\documentclass[uplatex]{jsarticle}
\usepackage{amsmath,amssymb,amsthm}

% --- theorem env ---
\theoremstyle{definition}
\newtheorem{definition}{定義}[section]
\newtheorem{theorem}{定理}[section]
\theoremstyle{remark}
\newtheorem{remark}{注意}[section]

\newcommand{\calF}{\mathcal{F}}
\newcommand{\R}{\mathbb{R}}

\begin{document}

\title{測度論：最小サンプル（定義＋基本定理）}
\author{}
\date{}
\maketitle

\section{定義}

\begin{definition}[$\sigma$-代数]
集合 $X$ の部分集合族 $\calF\subset 2^X$ が $X$ 上の $\sigma$-代数であるとは、
\begin{enumerate}
  \item $X\in \calF$、
  \item $A\in \calF \Rightarrow X\setminus A\in \calF$、
  \item 任意の列 $(A_n)_{n\in\mathbb{N}}\subset\calF$ に対し
        $\displaystyle \bigcup_{n=1}^{\infty}A_n\in\calF$
\end{enumerate}
が成り立つことをいう。
このとき $(X,\calF)$ を可測空間という。
\end{definition}

\begin{definition}[測度]
可測空間 $(X,\calF)$ 上の写像 $\mu:\calF\to[0,\infty]$ が測度であるとは、
\begin{enumerate}
  \item $\mu(\varnothing)=0$、
  \item （可算加法性）互いに素な可測集合列 $(A_n)_{n\in\mathbb{N}}\subset\calF$
        に対して
        \[
          \mu\!\left(\bigcup_{n=1}^{\infty}A_n\right)=\sum_{n=1}^{\infty}\mu(A_n)
        \]
\end{enumerate}
が成り立つことをいう。
このとき $(X,\calF,\mu)$ を測度空間という。
\end{definition}

\begin{definition}[可測関数]
$(X,\calF)$ を可測空間とする。
関数 $f:X\to\R$ が（$\calF$-）可測であるとは、
任意の実数 $a\in\R$ に対して
\[
  \{x\in X : f(x) > a\}\in\calF
\]
が成り立つことをいう。
\end{definition}

\section{基本定理：測度の連続性}

\begin{theorem}[下からの連続性]
測度空間 $(X,\calF,\mu)$ をとる。
可測集合列 $(A_n)_{n\in\mathbb{N}}\subset\calF$ が
\[
  A_1\subset A_2\subset \cdots
\]
を満たすとき、$A=\bigcup_{n=1}^\infty A_n$ とおけば
\[
  \mu(A)=\lim_{n\to\infty}\mu(A_n)
\]
が成り立つ。
\end{theorem}

\begin{proof}
$B_1:=A_1$、$B_n:=A_n\setminus A_{n-1}\ (n\ge2)$ と定める。
すると $(B_n)$ は互いに素で、また
\[
  A_n=\bigcup_{k=1}^n B_k,\qquad
  A=\bigcup_{k=1}^\infty B_k
\]
が成り立つ。測度の可算加法性より
\[
  \mu(A)=\sum_{k=1}^\infty \mu(B_k),\qquad
  \mu(A_n)=\sum_{k=1}^n \mu(B_k)
\]
である。よって $n\to\infty$ とすると
$\mu(A_n)\to \sum_{k=1}^\infty \mu(B_k)=\mu(A)$ が得られる。
\end{proof}

\begin{theorem}[上からの連続性]
測度空間 $(X,\calF,\mu)$ をとる。
可測集合列 $(A_n)_{n\in\mathbb{N}}\subset\calF$ が
\[
  A_1\supset A_2\supset \cdots
\]
を満たし、さらに $\mu(A_1)<\infty$ を仮定する。
$A=\bigcap_{n=1}^\infty A_n$ とおけば
\[
  \mu(A)=\lim_{n\to\infty}\mu(A_n)
\]
が成り立つ。
\end{theorem}

\begin{proof}
$C_n:=A_1\setminus A_n$ とおくと $(C_n)$ は増大列で、
\[
  \bigcup_{n=1}^\infty C_n = A_1\setminus \bigcap_{n=1}^\infty A_n
  =A_1\setminus A
\]
が成り立つ。下からの連続性より
\[
  \mu(A_1\setminus A)=\lim_{n\to\infty}\mu(C_n)
\]
である。一方 $\mu(A_1)<\infty$ より
\[
  \mu(C_n)=\mu(A_1\setminus A_n)=\mu(A_1)-\mu(A_n)
\]
が使えるので、
\[
  \mu(A_1)-\mu(A)=\lim_{n\to\infty}\bigl(\mu(A_1)-\mu(A_n)\bigr)
\]
となり、両辺から $\mu(A_1)$ を消して
$\mu(A)=\lim_{n\to\infty}\mu(A_n)$ を得る。
\end{proof}

\begin{remark}
上からの連続性では $\mu(A_1)<\infty$ の仮定が重要である。
これは「無限大の差 $\infty-\infty$」を避けるための条件と考えられる。
\end{remark}

\end{document}
